\bookchapter{Capítulo 17: Conexão do celular com outros dispositivos}

\chapterintro{O celular na sua mão é, na verdade, um super controle remoto e central de dados. Ele não só faz ligações, como também pode conectar computador, TV, relógio, caixa de som e carro, deixando de serem ilhas isoladas. Neste capítulo, vamos ver como fazer esses dispositivos “formarem um time”.}

\booksection{I. Integração com o computador: celular e computador não são mais dois mundos separados}

Sinceramente, o que mais me irritava antes era ter que achar o cabo de dados e esperar a transferência quando queria editar no computador uma foto tirada no celular. Depois descobri que, na verdade, nem precisa de cabo — uma boa integração é dez vezes mais rápida que o cabo.

\begin{itemize}
  \item \textbf{Ecossistema Apple (AirDrop + Área de Transferência Universal)}: Se você usa iPhone com Mac, o AirDrop é a primeira função que você deve aprender. No celular, toque em Compartilhar AirDrop selecione seu Mac. A foto chega instantaneamente, sem usar dados móveis nem Wi-Fi — usa Bluetooth e Wi-Fi Direct. Outra função pouco conhecida é a "Área de Transferência Universal" — copie um texto no iPhone e cole diretamente no Mac. É como mágica.
  \item \textbf{Android + Windows (Vincular ao Windows)}: A Microsoft tem um app chamado "Vincular ao Windows". Já vem instalado nos celulares Samsung; outros Androids é só baixar na loja de aplicativos. Depois de parear, você pode ver notificações, responder SMS, ver fotos e até fazer chamadas diretamente do computador. Uso há dois anos: em reuniões, não preciso ficar olhando o celular o tempo todo — uma notificação na tela do computador já resolve.
  \item \textbf{Colaboração multitela das marcas}: A "Colaboração Multitela" da Huawei e o "Xiaomi Share" da Xiaomi vão ainda mais longe — espelham diretamente a tela do celular no computador. Você pode arrastar arquivos do celular para o computador com o mouse, e também usar o TikTok e responder ao WeChat no computador. Cópia e cola entre dispositivos e arraste de arquivos, tudo funciona com uma fluidez impressionante.
\end{itemize}

\booksection{II. Espelhamento para TV: assistir em família}

No Ano Novo Chinês, quando fui visitar a família, minha prima queria mostrar na TV o vídeo da viagem de formatura que estava no celular dela. Depois de muito quebra-cabeça, descobrimos que era muito simples:

\begin{itemize}
  \item \textbf{iPhone TV}: Deslize para baixo do canto superior direito para abrir o Centro de Controle, toque no ícone de dois retângulos sobrepostos (Screen Mirroring), encontre o nome da sua TV e toque — conecta na hora. O pré-requisito é que a TV e o celular estejam na mesma rede Wi-Fi, e a TV suporte AirPlay (a maioria das smart TVs suporta).
  \item \textbf{Android TV}: Cada marca chama de um jeito. Huawei chama de "Projeção sem fio", Xiaomi de "Transmitir", OPPO/vivo de "Interação multitela". Geralmente está no painel de atalhos. Se não encontrar, procure por "transmitir" ou "espelhar" nas configurações.
  \item \textbf{Transmissão direta pelo app de vídeo}: Se você só quer assistir iQiyi, Bilibili, YouTube na TV, é ainda mais simples — na interface de reprodução, procure um ícone de retângulo com símbolo de Wi-Fi, toque nele, selecione a TV e o vídeo vai para lá. Enquanto isso, você pode fazer outras coisas no celular, pois o vídeo continua rodando na TV, sem interferir.
  \item \textbf{O que fazer se o espelhamento não funcionar?}90\% dos casos é porque não estão na mesma rede Wi-Fi. Os outros 10\% é porque a função de espelhamento da TV não está ativada — vá nas configurações da TV e procure por "Interação multitela" ou "AirPlay".
\end{itemize}

\booksection{III. Com smartwatch/pulseira inteligente: o monitor de saúde no seu pulso}

Comprei uma pulseira para minha mãe. No começo ela disse: “Nem faço exercício, pra que comprar isso?”. Duas semanas depois, ela descobriu que dormia apenas pouco mais de cinco horas por dia e começou a ajustar ativamente sua rotina. Esse é o valor mais prático do relógio/pulseira — transformar dados de saúde invisíveis em números visíveis.

\begin{itemize}
  \item \textbf{Filtro de notificações}: Quando o celular está vibrando sem parar na bolsa ou no bolso, basta levantar o pulso para ver — é uma mensagem do chefe ou uma promoção do app de compras? Em um segundo você decide. As mensagens que não valem tirar o celular do bolso, deixe-as para lá.
  \item \textbf{Sincronização de dados de saúde}: O relógio monitora 24 horas por dia: frequência cardíaca, oxigênio no sangue, sono, passos. Esses dados são automaticamente consolidados no app de Saúde do celular. Com o tempo, você vê suas tendências — este mês andou muito pouco? A qualidade do sono melhorou esta semana?
  \item \textbf{Sincronização de cartão de transporte/acesso}: Cartões de ônibus e de acesso cadastrados no celular são automaticamente sincronizados com o relógio na maioria das marcas. Ao pegar o metrô, nem precisa tirar o celular do bolso — é só encostar o pulso na catraca. Dá uma sensação real de "o futuro já chegou".
  \item \textbf{Ferramenta para encontrar o celular}: O relógio geralmente tem uma função "Encontrar celular". Um toque e o celular emite um som. Eu mesmo uso essa função muito mais do que esperava — sempre deixo o celular em algum lugar e esqueço onde.
\end{itemize}

\booksection{IV. Com casa inteligente: o celular é o “interruptor geral” da casa}

Casa inteligente parece algo complicado, mas se resume a três coisas: controle remoto pelo celular, configuração de regras de automação e controle por voz.

Vou contar alguns usos meus na prática, e você vai ver como isso é realmente útil:

\begin{itemize}
  \item \textbf{Ar-condicionado}: No verão, a caminho de casa, ligo o ar-condicionado antecipadamente pelo celular. Entro em casa e já está fresquinho. Com um plugue inteligente para ar-condicionado de cerca de cem reais, você consegue isso sem trocar o aparelho.
  \item \textbf{Luzes}: Deitado na cama antes de dormir, com preguiça de levantar para apagar a luz — é só tocar no celular ou dizer "Ei Siri / Xiao Ai, apague a luz".
  \item \textbf{Robô aspirador}: Antes de sair de casa, toque em "Iniciar limpeza" no celular. Quando voltar, o chão está limpo. Também dá para programar a limpeza diária.
  \item \textbf{Câmera}: Quem tem idosos, crianças ou pets em casa pode dar uma olhada na situação a qualquer momento pelo celular.
\end{itemize}

Qual ecossistema escolher? Resumo em uma frase:

\begin{itemize}
  \item Usa iPhone Priorize dispositivos compatíveis com Apple HomeKit (um pouco mais caros, mas com boa proteção de privacidade).
  \item Quer custo-benefício com muitos dispositivos O ecossistema Xiaomi Mijia é o mais completo: lâmpadas, plugues, robôs aspiradores, tem de tudo.
  \item Usa celular Huawei Priorize HarmonyOS Connect, a melhor experiência com dispositivos Huawei.
\end{itemize}

\booksection{V. Fones de ouvido Bluetooth e caixas de som: a liberdade sem fio}

Trocar o fone com fio pelo Bluetooth foi, para mim, a melhoria mais perceptível na experiência com o celular. Nunca mais preciso tirar um emaranhado de fios da bolsa, nem me preocupar com o fio enganchando na maçaneta.

\begin{itemize}
  \item \textbf{Pareamento é simples}: Ao ligar o fone pela primeira vez, geralmente ele entra automaticamente no modo de pareamento (luz piscando). Abra o Bluetooth do celular e ele aparecerá na busca. Toque e conecta. Depois disso, toda vez que tirar do estojo, conecta automaticamente, sem fazer nada.
  \item \textbf{Mesma marca tem vantagens}: AirPods com iPhone mostram automaticamente um pop-up com o nível da bateria e alternam perfeitamente entre vários dispositivos Apple. Huawei FreeBuds com Huawei e fones Xiaomi com celulares Xiaomi também têm pop-ups e otimizações semelhantes.
  \item \textbf{Caixa de som Bluetooth}: Em reuniões familiares e acampamentos ao ar livre, uma caixa de som Bluetooth dobra o clima. Depois de conectada, o celular é sua mesa de DJ.
  \item \textbf{Atenção à latência}: Ao jogar com fone Bluetooth, pode haver um leve atraso (o tiro soa meio segundo depois do disparo). Se a latência for um problema, use fone com fio ou compre um fone Bluetooth com suporte a "modo de baixa latência".
\end{itemize}

\begin{tcbinfo}{Tecnologia UWB (Ultra-Wideband) — sensação espacial mais precisa que o Bluetooth}
Além do Bluetooth, os celulares mais novos também trazem um chip UWB (Ultra-Wideband, banda ultralarga). Sua capacidade é a\textbf{medição precisa de distância e posicionamento espacial}— ele consegue calcular a que distância e em que direção está outro dispositivo, com precisão de centímetros. O Bluetooth só diz “está por perto”; o UWB diz “está a 1,2 metro à sua frente e à esquerda”. Aplicações comuns:

\begin{itemize}
  \item \textbf{AirTag / Busca precisa}: Quando o iPhone busca um AirTag e mostra a navegação por seta "7 metros, vá para a esquerda", é graças ao UWB. O Bluetooth comum só informa "está ou não está no alcance".
  \item \textbf{Chave digital do carro}: iPhones e flagships Android com UWB podem funcionar como chave do carro: destrava ao se aproximar e trava ao se afastar, além de diferenciar a posição do motorista e do passageiro, prevenindo ataques de retransmissão.
  \item \textbf{Navegação interna}: Em shoppings e aeroportos, os beacons UWB permitem navegação no nível de "corredor", muito mais precisa que o GPS.
  \item \textbf{Modelos compatíveis}: iPhone 11 e posteriores possuem chip UWB (U1/U2). No lado Android, alguns flagships da Xiaomi, Samsung, vivo e OPPO também já o possuem.
\end{itemize}

\end{tcbinfo}

\begin{tcbinfo}{LE Audio e Auracast — a nova geração de padrão de áudio Bluetooth}
O áudio Bluetooth clássico (A2DP), usado por duas décadas, está sendo substituído pelo\textbf{LE Audio}(Áudio de Baixa Energia), o novo padrão de áudio a partir do Bluetooth 5.2:

\begin{itemize}
  \item \textbf{Mais economia de energia, melhor qualidade de som}: Usa o novo codec LC3, com qualidade de som melhor e consumo menor que o SBC do padrão antigo. A autonomia dos fones é significativamente prolongada.
  \item \textbf{Áudio de dois canais}: Transmissão independente para os fones esquerdo e direito. A troca para modo mono não causa desconexão.
  \item \textbf{Auracast: áudio por broadcast}: A função carro-chefe do LE Audio. Um celular/TV pode transmitir áudio simultaneamente para\textbf{inúmeros}fones de ouvido — salas de embarque de aeroportos, museus, academias, salas de aula: todos podem usar uma única fonte de transmissão e várias pessoas ouvem o mesmo conteúdo com seus próprios fones.
  \item \textbf{Situação atual}: Os fones e celulares flagships lançados em 2024-2026 geralmente suportam LE Audio. Dispositivos antigos precisam de upgrade de hardware, não sendo possível desbloquear por software. Ao comprar fones novos, verifique se têm a indicação LE Audio/Auracast.
\end{itemize}

\end{tcbinfo}

\booksection{VI. Com sistema veicular: o parceiro de segurança ao dirigir}

Mexer no celular dirigindo é ilegal e pode ser fatal. Mas às vezes você realmente precisa de navegação, atender chamadas, tocar música. É aí que o CarPlay (Apple) e o Android Auto (Android) entram em cena.

\begin{itemize}
  \item \textbf{Como conectar}: Se seu carro for compatível (a maioria dos carros novos dos últimos anos é), conecte o celular com o cabo de dados na porta USB do carro. A tela do sistema veicular mudará automaticamente para a interface CarPlay ou Android Auto. Alguns carros novos suportam conexão sem fio, ainda mais prático.
  \item \textbf{O que dá para fazer}: Navegação (muito melhor que o GPS nativo do carro), chamadas (atender pelo volante), música (apps de música do celular direto no som do carro), enviar mensagens por voz (Siri ou assistente ajudam a ler e responder mensagens).
  \item \textbf{Carros antigos também podem usar}: Se seu carro for mais antigo e não tiver suporte, gaste algumas centenas de reais em uma tela CarPlay externa, fixe no painel, e a experiência é quase a mesma.
\end{itemize}

\booksection{VII. Segurança e privacidade na conexão entre dispositivos}

Quanto mais dispositivos conectados, maior a superfície de risco de vazamento de privacidade. Alguns pontos importantes:

\begin{itemize}
  \item Antes de espelhar a tela em público, verifique se há mensagens privadas na barra de notificações. Não deixe a família toda ver o que não deve.
  \item Antes de dar ou vender um dispositivo para outra pessoa, desvincule-o — vá nas configurações do celular, encontre o dispositivo e selecione "Remover" ou "Cancelar pareamento".
  \item Câmeras de segurança inteligentes: compre de marcas conhecidas, com criptografia. Não compre câmeras baratas de marcas desconhecidas. Já se ouviu notícias demais de câmeras hackeadas sendo transmitidas ao vivo.
\end{itemize}

\begin{tcbpractice}{Experimente: espelhe a tela do celular na TV e configure um cenário inteligente}
\begin{enumerate}
  \item Certifique-se de que o celular e a TV estão na mesma rede Wi-Fi.
  \item Na TV, encontre a função "Screen Mirroring" ou "AirPlay" e ative-a.
  \item No celular: iPhone — toque em «Screen Mirroring» no Centro de Controle; Android — deslize para baixo e procure «Projeção sem fio».
  \item Escolha sua foto favorita e espelhe na TV para a família ver.
  \item Se tiver dispositivos de casa inteligente, tente criar um "Modo Fora de Casa" no app: apagar luzes, desligar ar-condicionado e ligar o robô aspirador com um só toque.
\end{enumerate}

\end{tcbpractice}

\begin{tcbinfo}{Caso 1: Otimização prática da autonomia da bateria}
\begin{itemize}
  \item \textbf{Cenário}: Usuário com Xiaomi 14 com autonomia ruim, descarregando em meio dia.
  \item \textbf{Passos}: (1) Ajustes Bateria Ver ranking de consumo (2) Identificar apps anômalos em segundo plano (3) Desligar 5G e mudar para 4G (4) Restringir segundo plano (5) Ativar modo de economia de energia (6) Remover capa para dissipar calor ao carregar.
  \item \textbf{Resultado esperado}: Autonomia prolongada em 30-50\%, dados de saúde da bateria legíveis.
\end{itemize}

\end{tcbinfo}
